

En esta tarea se implementaron cuatro estrategias para resolver el problema de selección de capítulos por anime bajo restricciones de minutos y energía: fuerza bruta, programación dinámica y dos variantes greedy (greedy 2: satisfacción por energía; greedy 1: satisfacción por minuto). A partir de los experimentos realizados, se observó que la programación dinámica entrega la solución óptima en los casos evaluados, pero a costa de un mayor consumo de tiempo y memoria que los métodos greedy.

La fuerza bruta permitió verificar la correctitud de las otras implementaciones en instancias pequeñas, pero su costo computacional crece rápidamente, por lo que deja de ser práctica para tamaños medianos o grandes. Por su parte, los algoritmos greedy resultaron mucho más rápidos y livianos en memoria, aunque con una calidad de solución inferior a la óptima. Entre ambos, la variante basada en satisfacción por energía mostró un comportamiento más cercano al óptimo que la basada en satisfacción por minuto.

En síntesis, los resultados confirman el trade-off esperado entre exactitud y eficiencia. La programación dinámica es la alternativa adecuada cuando se requiere la mejor solución posible, mientras que las heurísticas greedy son útiles cuando se prioriza rapidez y bajo consumo de recursos. La fuerza bruta cumple un rol principalmente de validación en casos pequeños.
